\chapter{اللقاء التاسع عشر: بداية الجزء الثاني من سورة البقرة (سيقول السفهاء)}

\section{مقدمة: استئناف الدروس بعد مراجعة الجزء الأول}
السلام عليكم ورحمة الله وبركاته. الحمد لله رب العالمين، وأشهد أن لا إله إلا الله وحده لا شريك له، وأشهد أن محمداً عبده ورسوله. اللهم صل على محمد وعلى آل بيته وذريته كما صليت على آل إبراهيم إنك حميد مجيد، اللهم بارك على محمد وعلى آل بيته وذريته كما باركت على آل إبراهيم إنك حميد مجيد.

أهلاً ومرحباً بالطلاب الكرام، هذا هو الدرس السابع عشر (التاسع عشر بترقيم الدورة) من اجتماعنا على مدارسة كتاب تفسير الإمام \rawi{ابن جرير الطبري} عليه رحمة الله. وقد انقطعنا فترة بعد إنهاء الجزء الأول من هذا التفسير ليراجع الطلاب دروسهم ويُحضر كل واحد منهم ما وُكل إليه من المهام في استخراج منهج الطبري. وسنستأنف هذه الدروس لنركز على إنهاء هذا الكتاب، وسيكون معنا مسار آخر لقراءة الكتب التأسيسية في السنن (كموطأ مالك، ومصنف عبد الرزاق، ومصنف ابن أبي شيبة، وصحيح البخاري وغيرها) لأنها خير ما فُسر به القرآن، وهي تفسير تطبيقي عملي لكتاب الله.

نبدأ مستعينين بالله تبارك وتعالى في المجلد الثاني، من بداية الجزء الثاني من سورة البقرة.

\separator

\section{تأويل (سيقول السفهاء من الناس ما ولاهم عن قبلتهم)}
\isnad{قال الإمام الطبري رحمه الله:} «القول في تأويل قوله جل ثناؤه \textit{سَيَقُولُ السُّفَهَاءُ مِنَ النَّاسِ}... وهم اليهود وأهل النفاق. وإنما سماهم الله عز وجل سفهاء لأنهم سفهوا الحق، فتجاهلت أحبار اليهود وتعاظمت جهالهم... عن اتباع محمد ﷺ إذ كان من العرب ولم يكن من بني إسرائيل، وتحير المنافقون فتبلدوا».

وذكر \rawi{الطبري} اختلاف المفسرين فيهم؛ فمنهم من قال اليهود (وهو مروي عن \rawi{مجاهد} و\rawi{البراء} و\rawi{ابن عباس})، ومنهم من قال المنافقون (وهو مروي عن \rawi{السدي}). وجمع الطبري بينهما في المعنى.

\begin{fawaid}[طريقة القرآن في الإخبار والتوجيه (الخبر والإنشاء)]
يجب أن نعلم أن كتاب الله هو خبر وإنشاء (أمر ونهي)؛ فالله يُعلمك الخبر ليُمهد به للعمل. هنا أخبرنا الله بما سيقوله السفهاء من اليهود والمنافقين في تغيير القبلة ومكرهم، ثم علمنا كيف نتصرف ونجيبهم بقوله (قُل لِّلَّهِ الْمَشْرِقُ وَالْمَغْرِبُ). وهذا هو معنى صبغة الله والاستقامة، فالمؤمن الذي يتبع أمر الله هو الذي يصدق بالخبر ويعمل بالأمر، فيواجه الفتن والابتلاءات باليقين والعمل الصالح.
\end{fawaid}

\begin{fawaid}[الفرق بين أنواع الخلاف وأصنافه]
ذكر الطبري في هذه الآيات خلافات متنوعة. ومن المهم لطالب العلم التفريق بين:
1. \textbf{أنواع الخلاف:} (خلاف لفظي، اختلاف تنوع باختلاف الجهة، اختلاف تضاد).
2. \textbf{أصناف الخلاف:} وهي مجالات الخلاف التفسيري، كأن يختلفوا في: مدة الصلاة لبيت المقدس، سبب تحويل القبلة، سبب النزول، قائل القول، عود الضمير، أو الإعراب.
\end{fawaid}

\separator

\section{تأويل (قل لله المشرق والمغرب يهدي من يشاء إلى صراط مستقيم)}
\isnad{قال أبو جعفر:} «يعني جل ثناؤه بذلك: قل يا محمد لهؤلاء الذين قالوا لك ولأصحابك ما ولاكم عن قبلتكم من بيت المقدس... لله ملك المشرق والمغرب، يعني بذلك ملك ما بين قطري مشرق الشمس ومغربها... يهدي من يشاء من خلقه فيسدده ويوفقه إلى الطريق القويم».

\begin{fawaid}[التسليم المطلق لأحكام الله]
علمنا الله كيف نجيب من ينكر شيئاً من شريعتنا؛ وهو أن تعلم أننا عباد لله تبارك وتعالى، نسلم له بلا قيد ولا شرط. فحين سألوا مستنكرين تحويل القبلة، لم يجادلهم ببيان الحكمة، بل أجابهم بالتسليم: \textit{قُل لِّلَّهِ الْمَشْرِقُ وَالْمَغْرِبُ}. فالمؤمن يقبل أمر الله ويصدق خبره سواء علم الحكمة أم لم يعلمها، ومن توقف عن العمل حتى يعلم الحكمة أو أورد الشبهات والاعتراضات فهو يعبد الله على حرف.
\end{fawaid}

\separator

\section{تأويل (وكذلك جعلناكم أمة وسطا لتكونوا شهداء على الناس)}
\isnad{قال أبو جعفر:} «يعني جل ثناؤه بقوله \textit{وَكَذَٰلِكَ جَعَلْنَاكُمْ أُمَّةً وَسَطًا}... ففضلناكم بذلك على غيركم من أهل الأديان... وأما الوسط فإنه في كلام العرب يقال منه فلان واسط الحسب في قومه أي متوسط الحسب (إذا أرادوا بذلك الرفع من حسبه)... وأرى أن الله تبارك وتعالى إنما وصفهم بأنهم وسط لتوسطهم في الدين، فلا هم أهل غلو فيه غلو النصارى... ولا هم أهل تقصير فيه تقصير اليهود... ولكنهم أهل توسط واعتدال».

وساق الأخبار في أن الوسط هو العدل والخيار (كأحاديث \rawi{أبي سعيد الخدري} و\rawi{أبي هريرة} و\rawi{مجاهد} و\rawi{قتادة}). وأورد \rawi{الطبري} الأحاديث في شهادة أمة محمد ﷺ للأنبياء (كنوح عليه السلام) يوم القيامة على أممهم أنهم بلغوا رسالات ربهم.

\separator

\section{تأويل (وما جعلنا القبلة التي كنت عليها إلا لنعلم...)}
\isnad{قال أبو جعفر:} «يعني جل ثناؤه بقوله \textit{وَمَا جَعَلْنَا الْقِبْلَةَ الَّتِي كُنتَ عَلَيْهَا} ولم نجعل صرفك عن القبلة التي كنت على التوجه إليها يا محمد فصرفناك عنها \textit{إِلَّا لِنَعْلَمَ مَن يَتَّبِعُ الرَّسُولَ مِمَّن يَنقَلِبُ عَلَىٰ عَقِبَيْهِ}... فكان ذلك فتنة للناس وتمحيصاً للمؤمنين».

\begin{fawaid}[سنة الابتلاء بالتشريع والأحكام]
كما يبتلي الله عباده بالمصائب والخير والشر، فإنه يبتليهم بالتشريع وتغيير الأحكام (كتحويل القبلة من بيت المقدس إلى الكعبة)، ليميز الصادق المطيع من المعترض. فينقسم الناس حينها إلى مصدق محتسب، ومتردد شاك، ومنكر جاحد.
\end{fawaid}

\subsection{معنى (لنعلم) والرد على من استشكل العلم الحادث}
طرح \rawi{الطبري} إشكالاً متأثراً ببعض المتكلمين قائلاً: \isnad{فإن قال لنا قائل:} وما كان الله عالماً بمن يتبع الرسول ممن ينقلب على عقبيه إلا بعد اتباع المتبع... حتى قال \textit{إِلَّا لِنَعْلَمَ}؟

\isnad{فأجاب الطبري} بأن الله عالم بالأشياء كلها قبل كونها، ولكنه أوّل (لِنَعْلَمَ) بمعنى: «إلا ليعلم رسولي وحزبي وأوليائي»، مستدلاً بأن من شأن العرب إضافة فعل الأتباع للرئيس (كقولهم: فتح عمر العراق). وردّ قول من فسر (لنعلم) بـ (لنرى).

\begin{fawaid}[إثبات العلم الحادث لله بلا تأويل فاسد]
أخطأ الطبري ومن وافقه في تأويل (إِلَّا لِنَعْلَمَ) بصرف الضمير لغير الله (كجعله يعود على الرسول أو المؤمنين). والصواب المحكم أن العلم صنفان: علم بالشيء مقدراً قبل كونه (وهذا ثابت لله أزلاً)، وعلم بالشيء واقعاً حاصلاً. والله يحاسبنا على ما علمه منّا واقعاً وكسباً لا على مجرد علمه المسبق بنا. فالعلم الحادث (لنعلم واقعاً) لا ينفي العلم الأزلي المسبق، ولا حاجة لهذا التأويل المتكلف الذي لم يقله أحد من السلف، بل نبع من إشكالات المتكلمين (كدليل الأعراض) الذين نفوا الأفعال الاختيارية عن الله لئلا ينسبوا له الحدوث!
\end{fawaid}

\separator

\section{تأويل (وإن كانت لكبيرة إلا على الذين هدى الله)}
اختلف أهل التأويل في التي وصفها الله بأنها كانت كبيرة. فرأى بعضهم أنها (التولية والتحويلة)، ورأى آخرون أنها (القبلة بعينها)، ورأى آخرون أنها (الصلاة إلى القبلة الأولى).
\isnad{قال أبو جعفر مرجحاً:} «وأولى التأويلات عندي بالصواب: أن الكبيرة هي التولية والتحويلة (أي تحويل القبلة)... لأن القوم إنما كبر عليهم تحويل النبي ﷺ وجهه عن القبلة الأولى إلى الأخرى لا عين القبلة ولا الصلاة».

\begin{fawaid}[نعمة التوفيق للطاعة]
الهداية والتوفيق للطاعة فضل من الله يعين به عبده، فقوله: \textit{إِلَّا عَلَى الَّذِينَ هَدَى اللَّهُ} يبيّن أن من أُعين على طاعة (كقيام ليل أو حفظ قرآن) يجب أن يسند الفضل لله ويشكره ولا يغتر بعمله أو يزدري غيره، فالهداية نعمة وامتحان (ليبلوني أأشكر أم أكفر).
\end{fawaid}

\begin{fawaid}[الهداية نعمة وابتلاء معاً]
من المعاني التربوية الدقيقة هنا أن إعانة العبد على الطاعة ليست مجرد تشريف، بل هي أيضاً اختبار: هل يشكر هذه النعمة وينسبها إلى الله ويستعملها في مرضاته، أم يعجب بنفسه ويحتقر غيره؟ فمن فهم أن الهداية منحة من الله وامتحان له سلم من الكبر، وأقبل على الطاعة بخضوع وافتقار.
\end{fawaid}

\separator

\section{تأويل (وما كان الله ليضيع إيمانكم)}
ذكر \rawi{الطبري} الأخبار في سبب النزول أن المسلمين سألوا عمن مات من إخوانهم وهم يصلون إلى بيت المقدس قبل تحويل القبلة، هل بطلت أعمالهم؟ فنزلت: \textit{وَمَا كَانَ اللَّهُ لِيُضِيعَ إِيمَانَكُمْ}. أي: صلاتكم وتصديقكم.

\begin{fawaid}[الإيمان يطلق على الصلاة والأعمال]
في هذه الآية دلالة ظاهرة على إطلاق "الإيمان" على الصلاة والعمل. وهو ما يؤكد معتقد أهل السنة بأن الإيمان قول وعمل، وفيه رد صريح على المرجئة الذين أخرجوا العمل من مسمى الإيمان.
\end{fawaid}

\separator

\section{تأويل (قد نرى تقلب وجهك في السماء فلنولينك قبلة ترضاها...)}
\isnad{قال أبو جعفر:} «يعني بذلك جل ثناؤه: قد نرى يا محمد نحن تقلب وجهك في السماء... وإنما قيل ذلك له ﷺ فيما بلغنا بأنه كان قبل تحويل قبلته... يرفع بصره إلى السماء متوقعاً من الله جل ثناؤه أمره بالتحول نحو الكعبة... كراهة منه لقبلة اليهود».

ورجح الطبري أن (شطر) المسجد الحرام تعني: النحو والقصد والتلقاء. وناقش أقوال المفسرين في المكان الذي أُمر بتولية وجهه إليه؛ فرجح أنه "المسجد الحرام" نفسه (الكعبة ومحيطها)، وأن المولي وجهه شطر المسجد الحرام هو المصيب للقبلة وإن لم يحاذِ عيناً الكعبة ما دام متوجهاً لجهتها.

\separator

\section{تأويل (ولئن أتيت الذين أوتوا الكتاب بكل آية ما تبعوا قبلتك)}
\isnad{قال أبو جعفر:} «يعني جل ثناؤه... ولئن جئت يا محمد اليهود والنصارى بكل برهان وحجة... بأن الحق هو ما جئتهم به من فرض التحول... ما صدقوا به ولا تبعوا... وما أنت بتابع قبلتهم».

\begin{fawaid}[الاطلاع على بواطن الكفار لقطع التعلق بهم]
علّم الله نبيه ﷺ ما في قلوب أهل الكتاب من العناد لئلا يطمع في إرضائهم. ففي الآية دلالة بيّنة على نهي المؤمنين عن الركون إلى أعدائهم أو قبول نصيحتهم، لأن الله أطلعهم على ما يستبطنه الكفار من الحسد والبغض لئلا يُخدعوا بظاهرهم. وفيه تعزية للنبي ﷺ بأنه لا سبيل لإرضاء الفريقين معاً لتضادهما، فليلزم هدى الله الذي هو الحق.
\end{fawaid}

\separator

\section{تأويل (الذين آتيناهم الكتاب يعرفونه كما يعرفون أبناءهم)}
\isnad{قال أبو جعفر:} «يقول: يعرف هؤلاء الأحبار من اليهود والعلماء من النصارى أن البيت الحرام قبلتهم وقبلة إبراهيم والأنبياء... كما يعرفون أبناءهم».

\isnad{وأما قوله (وإن فريقا منهم ليكتمون الحق وهم يعلمون):} «وذلك الحق هو القبلة... وكتموا مع ذلك أمر محمد ﷺ وهم يجدونه مكتوباً عندهم... وهم يعلمون أنه ليس لهم كتمانه فيتعمدون معصية الله تبارك وتعالى».

\begin{fawaid}[الفرق بين المعرفة والعلم في القرآن]
المعرفة قد تكون مجرد الإدراك المجرد بلا عمل (يَعْرِفُونَهُ كَمَا يَعْرِفُونَ أَبْنَاءَهُمْ)، أما العلم في السياق القرآني فيتضمن الإدراك المتبوع بالخشية والعمل والاتباع، لذا عرف اليهود الحق لكنهم كتموه فاستحقوا الذم.
\end{fawaid}

\begin{fawaid}[أمانة العلم وعقوبة كتمانه]
العلم مسؤولية وميثاق. ومن رزقه الله العلم فكتمه أو حرفه أو استعمله للرياسة أو مظاهرة المجرمين فقد وقع في شعبة من النفاق وخصال اليهود. ولا قيمة لمجرد المعرفة إذا خلت من العمل والبذل لله.
\end{fawaid}

\separator

\section{تأويل (الحق من ربك فلا تكونن من الممترين)}
\isnad{قال الطبري رحمه الله:} «يعني جل ثناؤه: اعلم يا محمد أن الحق ما أعلمك ربك وأتاك من عنده... فلا تكونن من الشاكين في أن القبلة التي وجهتك نحوها هي الحق».

وطرح \rawi{الطبري} إشكالاً: أو كان النبي ﷺ شاكاً حتى يُنهى عن الشك؟
فأجاب: «ذلك من الكلام الذي يُخرج مخرج الأمر والنهي للمخاطب به والمراد به غيره (أي أمته)».

\begin{fawaid}[توجيه أوامر القرآن للنبي ﷺ]
تأويل الطبري بأن الخطاب مراد به غير النبي ﷺ (في مثل هذه الآية، وقوله: لئن أشركت ليحبطن عملك، واتق الله) هو تأويل محل نقد. فالصواب أن النبي ﷺ أول من يدخل في الخطاب، وتعليق الحكم على شرط لا يستلزم وقوعه. ونهيه عن الشك هو تثبيت له وموعظة، وهذا من كمال التشريع ولا منقصة فيه للنبي ﷺ، بل الأنبياء أشد الناس تكليفاً وأكملهم عبودية.
\end{fawaid}

\separator

\section{تأويل (ولكل وجهة هو موليها فاستبقوا الخيرات)}
\isnad{قال أبو جعفر:} «ولكل أهل ملة وجهة... فاستبقوا الخيرات: فبادروا وسارعوا... أي بادروا بالأعمال الصالحة شكراً لربكم وتزودوا في دنياكم لآخرتكم، فإني قد بينت لكم سبيل النجاة».

\separator

\section{تأويل (لئلا يكون للناس عليكم حجة إلا الذين ظلموا منهم فلا تخشوهم واخشوني)}
بين \rawi{الطبري} أن المراد بالناس الذين كان لهم حجة (خصومة) هم أهل الكتاب (اليهود)، وأن (الذين ظلموا منهم) هم مشركو قريش الذين قالوا "رجع إلى قبلتنا وسيرجع إلى ديننا".

\begin{fawaid}[إطلاق الحجة على الدعوى الباطلة]
لفظ "حجة" قد يُطلق على الخصومة والمجادلة والدعوى ولو كانت باطلة (كما في قوله: لئلا يكون للناس عليكم حجة)، فليس كل ما يُسمى حجة يكون صحيحاً وبرهاناً. فردّ الله على دعاوى المشركين بقوله: \textit{فَلَا تَخْشَوْهُمْ وَاخْشَوْنِي}، ليعلم المؤمن ألا يخاف من الشبهات والخصومات الباطلة التي تثار حول دينه.
\end{fawaid}

وردّ \rawi{الطبري} بقوة على اللغويين (كـ \rawi{أبي عبيدة} و\rawi{الفراء} و\rawi{الأخفش}) الذين زعموا أن (إلا) هنا بمعنى (الواو) للموالاة، أو بمعنى (لكن). وقرر أنه استثناء متصل صحيح، نافياً أن يكون للظالم حجة يعتد بها، وأن الإجماع يكفي لتخطئة قولهم.

\begin{fawaid}[تخطئة الأقوال اللغوية لمخالفتها الإجماع]
يرد الطبري تأويلات النحويين بشدة إذا خالفت إجماع المفسرين من الصحابة والتابعين وسياق الآيات، مقرراً أن "إجماع الحجة يكفي شاهداً على خطأ مقالتهم"، ولا يصح تفسير القرآن بالشواذ اللغوية التي تخالف المأثور.
\end{fawaid}

\separator

\section{تأويل (ولأتم نعمتي عليكم ولعلكم تهتدون * كما أرسلنا فيكم رسولا منكم...)}
\isnad{قال أبو جعفر:} «يعني تعالى ذكره بقوله \textit{وَلِأُتِمَّ نِعْمَتِي عَلَيْكُمْ}... فاكمل لكم به فضلي عليكم وأتم به شرائع ملتكم الحنيفية المسلمة... كما أرسلنا فيكم رسولاً منكم... يتلو عليكم آياتنا ويزكيكم ويعلمكم الكتاب والحكمة».

وبيّن أن (كما أرسلنا) متعلقة بـ (ولأتم نعمتي)، ورد على من جعلها متعلقة بـ (فاذكروني أذكركم) مدعياً التقديم والتأخير، معتمداً على اتساق المعنى وأفصح لغات العرب في عدم فصل المعاني.

\separator

\section{تأويل (فاذكروني أذكركم واشكروا لي ولا تكفرون)}
\isnad{قال أبو جعفر:} «يعني بذلك فاذكروني أيها المؤمنون بطاعتكم إياي فيما أمرتكم به وفيما نهيتكم عنه، أذكركم برحمتي ومغفرتي لكم... واشكروا لي أيها المؤمنون فيما أنعمت عليكم به من الإسلام والهداية... ولا تكفرون: ولا تجحدوا إحساني إليكم».

\begin{fawaid}[الشكر لا يكون تاماً إلا بالعمل]
فسر الطبري الشكر بأنه "الثناء على الرجل بأفعاله المحمودة". وهذا التفسير اللغوي فيه قصور عن المعنى الشرعي؛ فالشكر شرعاً لا يكتمل إلا بالعمل واستعمال النعم في طاعة المنعِم، كما قال تعالى: (اعملوا آل داود شكراً). والعبد الشكور هو العامل المطيع، لا المقتصر على مجرد الثناء والحمد باللسان.
\end{fawaid}

\separator

بهذا نكون قد أتممنا الجزء الأول من القرآن الكريم في تفسير الطبري بحمد الله. وسنستأنف في اللقاء القادم إن شاء الله. والسلام عليكم ورحمة الله وبركاته.
